\section{TEM 为什么“看起来像 Transformer”}

第二讲最后把 hippocampal associative memory 改写为 modern Hopfield network。现代 Hopfield update 与 softmax attention 共享指数相似度和 value 加权，因此 TEM 的 memory read 可以直接用 query、keys、values 表达。这个 reformulation 不只是一种类比，还改善 outer-product binding 的维度扩张，并允许额外 context 作为独立输入。

\subsection{从 TEM equation 到 Attention update}

第一张手写页把 TEM 的 recurrent/associative update 与 Transformer 公式并列。设当前 cue 为 $q_t$，历史 memory keys 为 $K=[k_1,\ldots,k_t]$，values 为 $V=[v_1,\ldots,v_t]$，则现代记忆读出为
$$
m_t=V\operatorname{softmax}(\beta K^{\top}q_t).
$$
其中 $q_t$ 是当前 hippocampal/structural cue，$k_\tau$ 是过去 memory address，$v_\tau$ 是要召回的 bound content。若只允许 $\tau\le t$，它也具备 causal attention 的时间约束。

\lecturefigure{slide-60-tem-transformer-equation.jpg}{把 TEM memory update 改写为 Transformer-like equation}{官方录像手写教学状态 V060；01:11:52}

\lecturefigure{slide-61-attention-query-retrieval.jpg}{Query、keys、values 视角下的 associative recall}{官方录像手写教学状态 V061；01:12:34}

\begin{importantbox}{Modern Hopfield 与 Attention 的共同计算}
当前 neural state 是 query，已存 memories 提供 keys 与 values，dot product 衡量相似度，softmax 选择并混合 memories。差别更多来自表示如何产生、memory 生命周期、训练目标与生物实现，而不是 read equation 的表面形式。
\end{importantbox}

\subsection{Relational memory 的更新步骤}

长手写页把流程拆成 key/query 生成、value encoding、causal mask、memory retrieval 与 recurrent state update。这里的 positional encoding 不是固定 sinusoid，而是 transition model 更新的 structural state；它带有环境 graph 的 inductive bias，因此一个 attention layer 就可能完成普通 Transformer 需要多层学习的关系推理。

\lecturefigure{slide-62-relational-memory-steps.jpg}{Relational memory 的 key/query/value 与 recurrent update}{官方录像手写教学状态 V062；01:14:56}

\begin{knowledgebox}{术语消化：positional encoding 在两种系统中不同}
标准 Transformer 常把 sequence index 编成位置；TEM 的 $g_t$ 更像由 action 累积得到的 latent graph position。两者都让相同 content 在不同结构位置产生不同表示，但 TEM 的位置更新由 transition dynamics 明确约束，归纳偏置更强。
\end{knowledgebox}

\subsection{Classic 与 modern Hopfield memory}

Neural-memory-softmax 页回顾：classic Hopfield 通过 recurrent energy descent 召回最接近 attractor，容量与干扰受限；modern Hopfield 使用指数相似度和显式 memory patterns，更新式近似 attention，并显著提高容量。课程利用的正是后者，而不是说所有 biological hippocampal recurrence 都使用 softmax。

\lecturefigure{slide-63-neural-memory-softmax.jpg}{Modern Hopfield memory 的 softmax retrieval}{官方录像手写教学状态 V063；01:15:22}

可用一个统一能量直觉表示：
$$
E(q)=-\frac{1}{\beta}\log\sum_{\mu}\exp(\beta k_\mu^{\top}q)
+\frac{1}{2}\|q\|^2.
$$
其中第一项鼓励 $q$ 接近某些 stored patterns，第二项约束状态大小。对能量做更新会得到 softmax-weighted memory combination。该式解释为何 memory capacity 与 retrieval sharpness 可以通过指数 kernel 改变。

\subsection{架构对应、缩放收益与限制}

Comparison 页把 LEC、MEC、hippocampus 与 content encoder、recurrent position state、attention memory 对齐；analysis 页继续比较 graph transitions 与 attention paths。Modern Hopfield 版本避免对 content 与 position 做巨大 outer product 再 flatten，额外 context 也可作为一个新输入流加入，复杂度从乘法式维度膨胀转向多路向量的近线性组合。

\lecturefigure{slide-64-tem-transformer-comparison.jpg}{TEM 与 Transformer-like components 的功能对应}{官方录像手写教学状态 V064；01:17:46}

\lecturefigure{slide-65-tem-analysis-diagrams.jpg}{对 TEM/Transformer correspondence 的结构分析}{官方录像手写教学状态 V065；01:18:20}

\teachervoice{Will 的结论非常克制：强 inductive bias 让任务更容易，因此一层 attention 就够；这不是证明“一层 Transformer 解决一切”。他把价值放在可解释性、神经预测、代码加速和额外 context 的可扩展绑定上。}

\begin{warningbox}{强归纳偏置带来的双面性}
若任务真由稳定 graph transition 和快速 content binding 构成，TEM 很高效；若结构本身频繁变化、动作不可观测或内容与结构强耦合，固定 factorization 可能成为错误偏置。与通用 Transformer 比较时必须匹配任务范围，而不是只比较层数。
\end{warningbox}

\subsection{结论与 Grid-cell Q\&A}

结论页总结双向收益：AI 的 modern memory 给 TEM 更好的 scaling 和实现；neuroscience 的 recurrent position code、content/structure separation 与 grid-like state 为 Transformer 设计和解释提供假设。最后 Q\&A 解释单个 grid cell 如何通过 electrode spike sorting 测得，不同 cells 可共享尺度/方向但相位平移，并形成多个 lattice-scale modules。

\lecturefigure{slide-66-conclusions.jpg}{第二讲结论：Neuroscience 与 AI 的双向关系}{官方录像手写教学状态 V066；01:19:42}

\teachervoice{学生问 grid template 是 evolution 还是 learning 的产物。Will 回答其规则结构在幼年动物中很早出现，显示强 developmental/evolutionary bias，但系统是否可被其他 relational tasks co-opt、灵活到什么程度仍不清楚。这个不确定性比一句“grid cells 是位置编码”更有教学价值。}

\subsection{本章小结}

Modern Hopfield retrieval 把 TEM memory 写成 attention：当前 state 是 query，历史 conjunctive memories 提供 keys/values，softmax 完成召回。Recurrent structural code 类似任务化 positional encoding。相似性产生可解释桥梁和实现收益，但 TEM 的强结构偏置、短任务范围与神经假设都限制了外推。

